\documentclass[10pt]{article}
\usepackage[margin=0.85in]{geometry}
\usepackage{amsmath,amssymb}
\usepackage[hidelinks]{hyperref}
\setlength{\parindent}{0pt}
\setlength{\parskip}{5pt}
\title{Selecting coupling modes during PCG: a diagnostic}
\author{}
\date{}

\begin{document}
\maketitle
\vspace{-2em}

Can the current residual guide a fixed-rank preconditioner update without
mixing error back into directions already resolved by preconditioned conjugate
gradients (PCG)? A global
Nystr\"om column replacement generally changes the preconditioned eigenvectors.
Scaled off-diagonal compression offers a different possibility: in an exact
two-block model, changing which singular modes are retained preserves a common
eigenvector basis. The hypothesis was that preserving this structure could make
adaptation useful in a hierarchical preconditioner. The experiment below tests
the selection rule in the exact two-block setting.

\textbf{Positive definite compression.}
Let $A=K+\lambda I$ with $K$ positive semidefinite and $\lambda>0$.
Let $A_{\mathrm{old}}$ be the symmetric positive definite (SPD) approximation
entering a compression level. Factor its diagonal blocks as
$(A_{\mathrm{old}})_{ii}=L_iL_i^\top$, and define
\[
  L=\operatorname{diag}(L_i),\qquad
  B=L^{-1}A_{\mathrm{old}}L^{-\top},\qquad
  P=\operatorname{diag}(V_iV_i^\top),\quad V_i^\top V_i=I.
\]
The scaling-and-compression construction described by Xing, Huang, and
Chow~\cite{xing2021} can be written as
\begin{equation}
  M=L\bigl[PBP+(I-P)\bigr]L^\top.                 \label{eq:spd}
\end{equation}
It retains the diagonal blocks and projects the scaled couplings. For nonzero
$w$, the middle matrix has quadratic form
$(Pw)^\top B(Pw)+\|(I-P)w\|_2^2>0$, so $M$ is SPD.
At successive levels, both the scaling and the compressed couplings must come
from the preceding SPD approximation. Replacing its couplings with blocks of
the original matrix does not inherit this argument.

\textbf{A common basis for changing the retained modes.}
For one partition of the original matrix, take exact diagonal-block factors.
Then
\[
 B=L^{-1}AL^{-\top}=
 \begin{bmatrix}I&T\\T^\top&I\end{bmatrix},\qquad
 T=U\Sigma V^\top,\qquad 0\leq\sigma_j<1.
\]
For a set $R$ of retained mode pairs, use
$P_R=\operatorname{diag}(U_RU_R^\top,V_RV_R^\top)$ in~\eqref{eq:spd}.
The resulting scaled matrix $B_R$ has eigenvectors
$w_j^{\pm}=2^{-1/2}(u_j,\pm v_j)^\top$.
Its eigenvalues are $1\pm\sigma_j$ for $j\in R$ and $1$ for omitted modes.
Thus every $B_R$ commutes with $B$: selection changes eigenvalues within the
same basis. In exact arithmetic, a residual component that is zero in this
basis remains zero after changing $R$ and restarting PCG. This does not preserve
the previous Krylov polynomial or establish a multistep convergence improvement.
Approximate bases and recursive compression can also destroy commutation.

\textbf{A residual-dependent selection criterion.}
At the current iterate, let $r=b-Ax$, $q=L^{-1}r$, and partition $q=(q_1,q_2)$.
Write $a_j=u_j^\top q_1$ and $b_j=v_j^\top q_2$. The contribution of an omitted
mode pair to the squared energy error in the inverse action is
\begin{equation}
 \eta_j=\frac{\sigma_j^2}{2}
 \left[\frac{(a_j+b_j)^2}{1+\sigma_j}
       +\frac{(a_j-b_j)^2}{1-\sigma_j}\right],\qquad
 \|M_R^{-1}r-A^{-1}r\|_A^2=\sum_{j\notin R}\eta_j, \label{eq:eta}
\end{equation}
where $\|z\|_A^2=z^\top Az$. This identity follows by resolving $q$ along
$w_j^+$ and $w_j^-$. At a single partition each mode pair has the same storage
cost, so retaining the $k$ largest $\eta_j$ minimizes this inverse-action error
over all size-$k$ subsets. The score is computable from $r$ and the coupling
SVD; it does not require the solution $A^{-1}b$.

\textbf{Diagnostic.}
The partition used a stable median split along the training feature with largest
range, exact diagonal-block factors, and $k\in\{16,32\}$ retained mode pairs.
The cases were synthetic smooth/random targets, Concrete, and Airfoil, with
seed 0 and $\lambda\in\{10^{-3},10^{-5}\}$. PCG began with the $k$ largest
singular values retained. Five branches compared uninterrupted static PCG,
unchanged-preconditioner restart at iteration 8, a complete subset replacement
by~\eqref{eq:eta} at iteration 8 followed by a frozen preconditioner, static
restarts at iterations 8/16/24, and subset replacements at those three iterations.
Branches shared the same prefix and used the directly checked original-system
tolerance $\|b-Ax\|_2/\|b\|_2\leq10^{-8}$. PCG ran in the common orthogonal
mode basis. Spectral products, original-matrix verification products, and dense
mappings for the original residual norm were counted separately.

\textbf{Observed outcome.}
All 80 solves reached the tolerance. One subset replacement required more
iterations than its matched static restart in all 16 comparisons; periodic
replacement lost 15 of 16. Its sole win was the smooth case at
$\lambda=10^{-3}$ and $k=32$: 28 versus 34 iterations. Yet all 64 adaptive
checkpoints reduced the exact squared inverse-action error by at least $80.2\%$.
An independent dense check gave relative error $6.96\times10^{-15}$ for
\eqref{eq:eta}, normalized commutator $1.97\times10^{-16}$, and a successful
Cholesky SPD check. The complete run used 4.88 CPU seconds; its raw record is
\texttt{results/two-block-oracle-diagnostic.json}.

Noncommutation is therefore unnecessary for this criterion to fail. An exactly
absent component stays absent, but a small component along a difficult mode is
not eliminated. Dropping that mode's correction can harm subsequent convergence
despite reducing current inverse-action error. These results retire selection
by~\eqref{eq:eta} and do not support proceeding to a hierarchical implementation.

\textbf{Resources and scope.}
The two exact diagonal factors require quadratic storage; the full
candidate SVD, residual-mapping caches, and dense kernel matrix are additional
resources, recorded separately from selection and solve work. Here $k$
counts singular mode pairs, not retained actual kernel columns. The experiment
therefore establishes neither an $O(nk)$ implementation nor an advantage at
the original column budget. For a hierarchy, node sizes, transfer matrices,
and candidate storage also prevent a fixed sum of ranks from being a fixed
memory budget.

Solution-dependent hierarchical refinement is already established: Bauer and
Bebendorf's block-adaptive cross approximation improves selected blocks using
the current approximate solution~\cite{bauer2019}. The present negative result
concerns replacing coupling modes at fixed capacity using~\eqref{eq:eta}; it
does not establish that all hierarchical adaptation fails. Any subsequent
mechanism still requires comparison with static allocation and total resource
accounting across the project's synthetic and real workloads.

\begin{thebibliography}{2}
\small
\bibitem{xing2021}
X. Xing, H. Huang, and E. Chow.
\emph{Efficient Construction of an HSS Preconditioner for Symmetric Positive
Definite $H^2$ Matrices}. SIAM J. Matrix Anal. Appl., 2021.
\href{https://doi.org/10.1137/20M1365776}{doi:10.1137/20M1365776}.

\bibitem{bauer2019}
M. Bauer and M. Bebendorf.
\emph{Block-adaptive Cross Approximation of Discrete Integral Operators}, 2019.
\href{https://arxiv.org/abs/1905.01918}{arXiv:1905.01918}.
\end{thebibliography}
\end{document}
